% =============================================================================
% SYSTEM INSTRUCTION TO LLM:
% This file is a PROMPT TEMPLATE.
% You MUST generate a FULL, FINAL, COMPILABLE LaTeX presentation
% based on this template.
%
% CRITICAL RULES:
% 1. DO NOT change anything in the Preamble (except the outline titles inside
%    \deckoutline, which you MUST edit).
% 2. You MUST delete all [STYLE REF] frames and the AI INSTRUCTION frame.
% 3. You MUST replace all placeholder texts with real content.
% 4. You MUST preserve the collegeBeamer theme and visual style.
% 5. PAGE BUDGET: total pages = 3 fixed pages (title + outline + Q&A)
%    + 1 transition page per section + your content pages.
%    Target 12--18 pages total. Keep each content frame to roughly
%    8--10 lines of body text.
% 6. Use \includegraphics only for the user image files listed in the request
%    (figures/...); when none are listed, no external image files exist. Draw every
%    other chart, diagram and table with tikz / pgfplots / booktabs.
% 7. Only these packages are available (already loaded in the preamble): collegebeamer,
%    tikz, pgfplots, amsmath, amssymb, pifont, booktabs, colortbl, multirow.
%    Do NOT add \usepackage lines.
% 8. OUTLINE: edit the section titles ONCE inside \deckoutline. Then write
%    \decktoc{0} for the outline page and \decktoc{k} right after \section{...}
%    of section k. NEVER copy outline text by hand onto multiple pages.
% 9. Block environments (block / alertblock / exampleblock) and lstlisting code
%    blocks are styled by the theme and are allowed on content pages.
% 10. Zebra-stripe a table with \rowcolors{2}{maincolor!7}{white} ONLY when it
%    has 4 or more data rows; leave short tables unstriped.
% 11. TABLES: give every table a solid navy header row -- start the header row
%    with \deckheadrow and use {\color{white}\bfseries ...} header cells, NO
%    \toprule above it -- exactly as in the [STYLE REF] table page. Keep
%    \midrule and \bottomrule.
% 12. CHART COLORS: first series maincolor (navy), second series airforceblue
%    (steel). Never introduce other hues. Every axis takes the [deck, ...] style.
% 13. The navy frametitle band shows the section tag automatically on its right
%    side; NEVER repeat the section name inside frame titles or body text.
% =============================================================================

\documentclass[aspectratio=169,10pt,xcolor=table]{beamer}

% -----------------------------------------------------------------------------
% 1. Preamble & Theme (DO NOT CHANGE)
% -----------------------------------------------------------------------------
\usepackage[zh,bjtu]{collegebeamer}

\usepackage{tikz}
\usetikzlibrary{positioning}
\usepackage{pgfplots}
\pgfplotsset{compat=1.17}
\usepackage{amsmath,amssymb}
\usepackage{pifont}
\usepackage{booktabs}
\usepackage{colortbl}
\usepackage{multirow}

% Disable collegeBeamer auto TOC insertion (\decktoc pages replace it)
\AtBeginSection[]{}

% Shared chart style: horizontal grid only (no vertical lines through bars),
% no axis arrow tips, readable label sizes, and a very light navy plot-area
% fill so charts read like report figures. Use [deck, ...] on every axis.
% Series colors are ALWAYS navy first, steel second: maincolor, airforceblue.
\pgfplotsset{
  deck/.style={
    axis y line=left,
    axis x line=bottom,
    axis line style={-, maincolor!60},
    axis background/.style={fill=maincolor!3},
    ymajorgrids=true,
    grid style={maincolor!12},
    tick label style={font=\scriptsize},
    label style={font=\scriptsize},
    legend style={font=\scriptsize, draw=none, fill=none},
  },
}

% Navy table header row. Start every table's header row with \deckheadrow
% (instead of a bare \rowcolor) -- the widened overhang closes the hairline
% gaps between the per-cell color panels.
\newcommand{\deckheadrow}{\rowcolor{maincolor}[\dimexpr\tabcolsep+0.4pt\relax][\dimexpr\tabcolsep+0.4pt\relax]}

% -----------------------------------------------------------------------------
% 1.1 Outline machinery (DO NOT CHANGE this block, EDIT titles in \deckoutline)
% -----------------------------------------------------------------------------
% \decktoc{0}  -> outline page: report-style agenda rows separated by hairlines
% \decktoc{k}  -> section divider: left third solid navy panel with the big
%                 white section number + name; right two thirds carry the
%                 agenda with entry k bold navy and all other entries dimmed
\newcounter{decktocurrent}
\newcommand{\decktocitem}[2]{%
  \ifnum\value{decktocurrent}=#1
    {\color{maincolor}\bfseries\makebox[2.2em][l]{0#1}#2\hfill\raisebox{0.25ex}{\scriptsize$\blacksquare$}}%
  \else
    \ifnum\value{decktocurrent}=0
      {\color{airforceblue}\makebox[2.2em][l]{0#1}}{\color{reporttext}#2}%
    \else
      {\color{reportmuted}\makebox[2.2em][l]{0#1}#2}%
    \fi
  \fi
  \par\vspace{0.55em}%
  {\color{reportline}\hrule height 0.4pt}%
  \vspace{0.80em}%
}
\newcommand{\decktoc}[1]{%
  \setcounter{decktocurrent}{#1}%
  \ifnum#1=0
    \begin{frame}{Outline}
      \large
      \vspace{0.8em}
      \deckoutline
    \end{frame}
  \else
    \begin{frame}[plain]
      \begin{tikzpicture}[remember picture, overlay]
        % 2pt bleed on the outer edges so no white hairline survives rasterizing
        \fill[maincolor] ([xshift=-2pt, yshift=2pt]current page.north west)
          rectangle ([xshift=0.34\paperwidth, yshift=-2pt]current page.south west);
        \node[anchor=north west, inner sep=0pt, text=white]
          at ([xshift=0.9cm, yshift=-1.7cm]current page.north west)
          {\fontsize{44}{44}\selectfont\bfseries 0#1};
        \draw[bandsteel, line width=0.6pt]
          ([xshift=0.9cm, yshift=-3.55cm]current page.north west) -- ++(1.15cm,0);
        \node[anchor=north west, inner sep=0pt, text=white, font=\large\bfseries,
              text width=\dimexpr0.34\paperwidth-1.8cm\relax]
          at ([xshift=0.9cm, yshift=-3.95cm]current page.north west) {\insertsection};
      \end{tikzpicture}%
      \hspace*{\dimexpr0.34\paperwidth+0.1cm\relax}%
      \begin{minipage}{\dimexpr0.66\paperwidth-2.1cm\relax}
        {\footnotesize\color{airforceblue}Agenda}\\[0.4em]
        {\color{reportline}\rule{\linewidth}{0.4pt}}
        \vspace{0.9em}
        \deckoutline
      \end{minipage}
    \end{frame}
  \fi
}

% LLM MUST EDIT: define each section title here ONCE (one \decktocitem per
% section, numbered 1..N). Every outline/transition page renders from this.
\newcommand{\deckoutline}{%
  \decktocitem{1}{Title of Part One}%
  \decktocitem{2}{Title of Part Two}%
  \decktocitem{3}{Title of Part Three}%
  \decktocitem{4}{Title of Part Four}%
  \decktocitem{5}{Summary \& Q\&A}%
}

% -----------------------------------------------------------------------------
% 1.2 Title page (DO NOT CHANGE): report cover. Uppercase letterspaced steel
%     kicker (institute, when provided), 27pt bold title, thin steel rule,
%     italic subtitle, author/date stack. Empty fields are skipped, no
%     floating gaps; the block sits in the calm left zone of the ring art.
% -----------------------------------------------------------------------------
\makeatletter
\setbeamertemplate{title page}{
    \hspace{-12mm}
    \begin{beamercolorbox}[wd=0.78\textwidth, sep=10pt, leftskip=8mm]{title}
        \setbox\@tempboxa=\hbox{\insertinstitute}%
        \ifdim\wd\@tempboxa>1pt
          % kicker reads the raw institute text (\beamer@shortinstitute):
          % \MakeUppercase chokes on \insertinstitute's embedded \def prefix
          {\footnotesize\color{airforceblue}\addfontfeatures{LetterSpace=16}%
           \def\inst##1{}\def\and{\space}%
           \MakeUppercase{\beamer@shortinstitute}}\\[11pt]
        \else
          {\color{airforceblue}\rule{24pt}{2pt}}\\[10pt]
        \fi
        {\fontsize{27}{32}\selectfont\bfseries\color{maincolor}\inserttitle}\\[9pt]
        {\color{airforceblue}\rule{0.32\textwidth}{0.5pt}}%
        \setbox\@tempboxa=\hbox{\insertsubtitle}%
        \ifdim\wd\@tempboxa>1pt
          \\[7pt]{\usebeamerfont{subtitle}\itshape\textcolor{reportmuted}{\insertsubtitle}}%
        \fi
        \setbox\@tempboxa=\hbox{\insertauthor}%
        \ifdim\wd\@tempboxa>1pt
          \\[10pt]{\usebeamerfont{author}\usebeamercolor[fg]{author}\insertauthor}%
        \fi
        \\[6pt]{\usebeamerfont{date}\usebeamercolor[fg]{date}\insertdate}
    \end{beamercolorbox}
}
\makeatother

% -----------------------------------------------------------------------------
% 2. Metadata (LLM MUST FILL)
% -----------------------------------------------------------------------------
\title[Short Title]{Full Presentation Title}
% Caller fills these three fields. Leave empty when unknown; never invent placeholders.
\subtitle{}
\author{}
\institute{}
\date{\today}

\begin{document}

% =============================================================================
% FRAME 1: TITLE PAGE
% =============================================================================
\maketitle

% =============================================================================
% FRAME 2: OUTLINE PAGE (rendered from \deckoutline; never hand-copy titles)
% =============================================================================
\decktoc{0}

% =============================================================================
% SECTION 1 (EXAMPLE)
% Pattern for EVERY section k: \section{...} then \decktoc{k}, then content
% frames. The transition page highlights the current section automatically.
% =============================================================================
\section{Title of Part One}

% FRAME 3: SECTION TRANSITION (EXAMPLE: SECTION 1)
\decktoc{1}

% =============================================================================
% CONTENT FRAME TEMPLATE (LLM MUST USE THIS STRUCTURE)
% =============================================================================
% Content page example 1: Text-left, figure-right layout
\begin{frame}{Content Example: Text Left, Figure Right}
\small
\begin{columns}[T]
\column{0.55\textwidth}
Write the core analysis content on the left side here.
This area is used for background description, problem definition,
logical reasoning, or methodology explanation.
Text should be concise and conclusion-oriented; avoid vague descriptions.

\vspace{0.6em}
\textbf{Key Takeaway:}
State the central message of this slide in one sentence.

\column{0.40\textwidth}
\centering
% Side panel: FEW categories only (4--6 short labels). Compact width. No xlabel.
% For 12-month / many-label series use the full-width chart on the next example page.
\begin{tikzpicture}
\begin{axis}[
    deck,
    width=5.4cm,
    height=3.6cm,
    ybar,
    bar width=13pt,
    enlarge x limits=0.15,
    ymin=0, ymax=100,
    symbolic x coords={Q1,Q2,Q3,Q4},
    xtick=data,
    ylabel={Sample index},
]
\addplot[fill=maincolor, draw=none] coordinates {
    (Q1,42) (Q2,55) (Q3,61) (Q4,78)
};
\end{axis}
\end{tikzpicture}
\end{columns}
\end{frame}

% Content page example 2: Text-top, figure-bottom layout
\begin{frame}{Content Example: Text Top, Figure Bottom}
\small

% Upper text area
When a series has many categories (for example 12 months), use this
full-width layout below -- not the narrow side-panel chart.
Keep the text concise so the wide chart remains readable.

\vspace{1em}

% Lower figure: MANY categories (12 months) → MUST be wide + rotated x labels
\centering
\begin{tikzpicture}
\begin{axis}[
    deck,
    width=11.5cm, height=4.2cm,
    enlarge x limits=0.04,
    x tick label style={rotate=45, anchor=east, font=\scriptsize},
    xlabel style={yshift=-2pt},
    ymin=1900, ymax=2900,
    symbolic x coords={Jan,Feb,Mar,Apr,May,Jun,Jul,Aug,Sep,Oct,Nov,Dec},
    xtick=data,
    ylabel={Sample level},
    xlabel={Month (illustrative)},
    legend style={at={(0.02,0.98)}, anchor=north west},
]
\addplot[
    thick,
    mark=*,
    maincolor,
    mark options={fill=white, draw=maincolor, scale=0.7},
] coordinates {
    (Jan,2040) (Feb,2085) (Mar,2160) (Apr,2330) (May,2345) (Jun,2320)
    (Jul,2420) (Aug,2505) (Sep,2630) (Oct,2740) (Nov,2655) (Dec,2620)
};
\addplot[
    thick,
    dashed,
    airforceblue,
    mark=square*,
    mark options={fill=white, draw=airforceblue, scale=0.55},
] coordinates {
    (Jan,2010) (Feb,2050) (Mar,2120) (Apr,2280) (May,2300) (Jun,2290)
    (Jul,2380) (Aug,2460) (Sep,2580) (Oct,2700) (Nov,2620) (Dec,2590)
};
\legend{Series A (sample), Series B (sample)}
\end{axis}
\end{tikzpicture}

\vspace{0.3em}
\centerstate{
    \tiny \emph{Source: data compiled from the source document.}
}
\end{frame}


% Content page example 3: Figure-left, text-right layout
\begin{frame}{Content Example: Figure Left, Text Right}
\small
\begin{columns}[c] % [c] = vertically centered; use [T] for top alignment if text is long

% Left: real TikZ structure diagram (layered pipeline)
\column{0.45\textwidth}
    \centering
    \begin{tikzpicture}[
        font=\scriptsize,
        box/.style={
            draw=maincolor, fill=maincolor!12, rounded corners=2pt,
            minimum height=0.62cm, text width=1.35cm, align=center
        },
        steelbox/.style={
            draw=airforceblue, fill=airforceblue!14, rounded corners=2pt,
            minimum height=0.62cm, text width=1.35cm, align=center
        },
        layer title/.style={font=\tiny\bfseries, text=maincolor}
    ]
        % Data layer (navy family)
        \fill[maincolor!6, rounded corners=3pt] (-0.25,1.05) rectangle (4.85,2.45);
        \node[layer title, anchor=north west] at (-0.15,2.40) {Data layer};
        \node[box] (s1) at (0.7,1.55) {Sensors};
        \node[box] (s2) at (2.5,1.55) {Ingest};
        \node[box] (s3) at (4.2,1.55) {Store};
        % Decision layer (steel family -- two-tone report look)
        \fill[airforceblue!8, rounded corners=3pt] (-0.25,-0.95) rectangle (4.85,0.70);
        \node[layer title, anchor=north west, text=airforceblue] at (-0.15,0.65) {Decision layer};
        \node[steelbox] (p1) at (0.7,-0.15) {Feature};
        \node[steelbox] (p2) at (2.5,-0.15) {Model};
        \node[steelbox] (p3) at (4.2,-0.15) {Act};
        % Edges
        \draw[->, thick, maincolor] (s1) -- (s2);
        \draw[->, thick, maincolor] (s2) -- (s3);
        \draw[->, thick, maincolor] (s2) -- (p1);
        \draw[->, thick, maincolor] (p1) -- (p2);
        \draw[->, thick, maincolor] (p2) -- (p3);
        \draw[->, thick, maincolor, dashed]
            (p3.south) -- ++(0,-0.28) -| (p1.south)
            node[pos=0.25, below, font=\tiny, text=maincolor] {feedback};
    \end{tikzpicture}

% Right: text analysis area (50\% width)
\column{0.50\textwidth}
    \textbf{Figure / Object Analysis:}

    This section provides a detailed interpretation of the visual element
    on the left. Suitable for presenting:
    \begin{itemize}
        \item Physical structure of experimental apparatus
        \item Specific module details of an algorithm or model
        \item Domain-specific data analysis
    \end{itemize}

    \vspace{1em}
    \textbf{Key Findings:}
    \begin{enumerate}
        \item \ctextbf{Observation 1}: The figure on the left shows\ldots
        \item \ctextbf{Observation 2}: Compared to expectations, the results exhibit\ldots
    \end{enumerate}
\end{columns}
\end{frame}

% =============================================================================
% [STYLE REF] TABLE PAGE (DELETE IN FINAL OUTPUT)
% =============================================================================
\begin{frame}{[STYLE REF] Table Style Example}
\small
\centering
\renewcommand{\arraystretch}{1.25}
\footnotesize
% BOARDROOM TABLE RULE: header row = solid navy with white bold text.
% Start the header row with \deckheadrow (NOT \toprule, NOT bare \rowcolor)
% and write each header cell as {\color{white}\bfseries ...}.
% In p{...} columns wrap the header cell in \multicolumn{1}{l}{...} so its
% baseline matches the other header cells.
% Numeric columns right-aligned (r); prose columns left (l / p{...}).
% Body keeps booktabs feel: light \midrule under header, \bottomrule to close.
% Add \rowcolors{2}{maincolor!7}{white} ONLY for tables with 4+ data rows.
\begin{tabular}{p{0.2\textwidth} r r l}
\deckheadrow
\multicolumn{1}{l}{\color{white}\bfseries Metric} & {\color{white}\bfseries A}
  & {\color{white}\bfseries B} & {\color{white}\bfseries Conclusion} \\
\midrule
Item 1 & 12 & 15 & Room for improvement \\
Item 2 & 20 & 18 & Stable \\
\bottomrule
\end{tabular}
\end{frame}

% =============================================================================
% [STYLE REF] TIKZ FLOW PAGE (DELETE IN FINAL OUTPUT)
% =============================================================================
\begin{frame}{[STYLE REF] Flowchart Style Example}
\centering
% Main path in navy; secondary / feedback nodes in steel (airforceblue)
\begin{tikzpicture}[
node distance=0.9cm and 0.7cm,
process/.style={rectangle,draw=maincolor,fill=maincolor!10,
minimum height=0.95cm,text width=1.9cm,align=center},
sideprocess/.style={rectangle,draw=airforceblue,fill=airforceblue!12,
minimum height=0.95cm,text width=1.9cm,align=center}
]
\node[process] (a) {Collect};
\node[process,right=of a] (b) {Analyze};
\node[process,right=of b] (c) {Decide};
\node[sideprocess,below=of b] (d) {Revise};
\draw[->,thick,maincolor] (a)--(b);
\draw[->,thick,maincolor] (b)--(c);
\draw[->,thick,airforceblue] (b)--(d)
    node[midway,right,font=\tiny,text=airforceblue] {gap};
\draw[->,thick,airforceblue,dashed] (d.west) -| (a.south)
    node[near end,left,font=\tiny,text=airforceblue] {retry};
\end{tikzpicture}
\end{frame}

% =============================================================================
% [STYLE REF] TWO-COLUMN ANALYSIS (DELETE IN FINAL OUTPUT)
% =============================================================================
\begin{frame}{[STYLE REF] Two-Column Analysis Example}
\small
\begin{columns}
\column{0.52\textwidth}
Left column for argumentation, definitions, and derivations.

\column{0.43\textwidth}
Right column for action items, illustrations, or comparisons.
\end{columns}
\end{frame}

% =============================================================================
% AI INSTRUCTION PAGE (DELETE IN FINAL OUTPUT)
% =============================================================================
\begin{frame}[allowframebreaks]{AI GENERATION INSTRUCTIONS}
You MUST:
\begin{itemize}
\item Delete all [STYLE REF] frames and this instruction frame
\item Replace all placeholders with real content
\item Edit the outline titles ONCE inside \texttt{\textbackslash deckoutline};
      use \texttt{\textbackslash decktoc\{0\}} for the outline page and
      \texttt{\textbackslash decktoc\{k\}} right after \texttt{\textbackslash section\{...\}}
      of section k -- never hand-copy the outline.
      \texttt{\textbackslash decktoc\{k\}} renders the navy left-panel section
      divider automatically; never build section pages by hand
\item Define \texttt{\textbackslash section\{...\}} before the first content frame of that section
\item Frame titles live in the navy header band; the section tag appears on the
      band's right automatically -- keep titles short and never repeat the
      section name in them
\item Use columns layout for all content pages
\item Keep information density high (no empty slides)
\item Respect the page budget: 3 fixed pages + 1 transition per section + content, 12--18 total
\item Give every table a solid navy header row: start the header row with
      \texttt{\textbackslash deckheadrow} + white bold cells, no
      \texttt{\textbackslash toprule} -- as shown on the [STYLE REF] table page
\item Style every chart with the \texttt{[deck, ...]} axis options; series
      colors are navy then steel: \texttt{maincolor}, \texttt{airforceblue}
\item Blocks (block / alertblock / exampleblock) and \texttt{lstlisting} are available
\item Zebra-stripe tables only when they have 4 or more data rows
\item Output FULL LaTeX code only
\end{itemize}
\end{frame}

% =============================================================================
% FINAL PAGE
% =============================================================================
\QApage

\end{document}
